O Robot Operating System (ROS) é um dos frameworks mais difundido para implementação de sistemas robóticos e ciber-físicos~\cite{ros1}. Utilizado amplamente em ambientes acadêmicos e industriais, o ROS movimenta um mercado estimado em 47,8 bilhões de dólares até 2025.\footnote{\href{https://www.thebusinessresearchcompany.com/market-insights/ros-based-robot-market-insights-2025}{https://www.thebusinessresearchcompany.com/market-insights/ros-based-robot-market-insights-2025}}
 Entretanto, com o fim do suporte ao ROS 1 em 2025,\footnote{\href{https://www.ros.org/blog/noetic-eol/}{https://www.ros.org/blog/noetic-eol/}}
 diversos sistemas correm risco de obsolescência. A migração para o ROS 2~\cite{ros2}, que incorpora avanços em segurança, gerenciamento de qualidade de serviço (QoS) e ciclo de vida dos nós é necessário mas carece de suporte estável ou documentação consolidada ou uma sistematização coerente para a transição.

Apesar de esforços recentes que aplicam princípios de Engenharia de Software ao projeto~\cite{malavolta2020} e à verificação e validação~\cite{caldas2024runtime,santos2021high} de sistemas ROS, a migração de sistemas legados envolve desafios relacionados à compreensão dos comportamentos existentes, à identificação de suas realizações na implementação e à avaliação de sua preservação na plataforma de destino. Assim, a migração de sistemas legados permanece um processo custoso, sujeito a erros e sem garantias de preservação das propriedades críticas originais.

Para investigar esse desafio, adotamos como estudo de caso a Body Sensor Network (BSN)~\cite{Gil2021BodySN}, um sistema auto-adaptativo de monitoramento médico desenvolvido pela equipe do LADECIC. Originalmente implementada em ROS 1, a BSN integra sensores como termômetro, oxímetro e eletrocardiograma em um bodyhub central, classificando o estado de saúde do paciente em níveis de risco. Sendo um sistema crítico, sua modernização deve preservar confiabilidade e adaptabilidade, ao mesmo tempo em que explora os benefícios introduzidos pelo ROS 2.

Este trabalho propõe uma abordagem sistemática para a migração de sistemas robóticos orientada pelo estabelecimento de uma \textit{Ground Truth} comportamental do sistema legado. O objetivo geral é propor e examinar essa abordagem por meio de um estudo de caso parcialmente retrospectivo da migração da BSN do ROS 1 para o ROS 2, investigando como a referência comportamental recuperada pode apoiar a análise das transformações realizadas e a avaliação da preservação dos comportamentos selecionados.

A abordagem é organizada a partir da fase de \textit{Ground Truth} Comportamental que contém a (i)obtenção de conhecimento sobre o sistema legado, (ii)a especificação de requisitos comportamentais por meio de cenários \textit{Behavior-Driven Development} (BDD) e o (iii)estabelecimento de uma \textit{Ground Truth} comportamental a partir da execução desses cenários no ROS 1. A fase de Migração utiliza esse \textit{Ground Truth} para orientar a (i) identificação das realizações dos comportamentos na implementação legada, o (ii)planejamento das transformações, a (iii)implementação no ambiente de destino e a (iv)validação comportamental no ROS 2. A integração dessas atividades busca relacionar os comportamentos especificados às responsabilidades de implementação e às transformações realizadas, considerando tanto as características da plataforma de origem quanto as do ambiente de destino.

A abordagem é examinada por meio de um estudo de caso empírico e retrospectivo, fundamentado na análise de documentação, código-fonte, cenários BDD, registros de execução e artefatos da migração da BSN. A avaliação é orientada por três questões de pesquisa: 
\begin{itemize}
    \item \textbf{RQ1 -- Aplicabilidade e limitações:} Como a abordagem proposta viabiliza a migração de sistemas robóticos? 
    \item \textbf{RQ2 -- Preservação comportamental:} Quais evidências recuperadas permitem avaliar a preservação dos cenários BDD selecionados durante a migração?
    \item \textbf{RQ3 -- Caracterização das transformações de migração:} Como as responsabilidades de implementação associadas aos contratos comportamentais selecionados foram transformadas durante a migração da BSN do ROS 1 para o ROS 2?    
\end{itemize}
 A investigação considera a exequibilidade das atividades propostas, a utilidade dos artefatos recuperados ou produzidos, as evidências de preservação comportamental e as transformações identificadas nas implementações.


Como resultados, o estudo apresenta a sistematização parcialmente retrospectiva do processo de migração da BSN, relacionando os contratos comportamentais selecionados às entidades e responsabilidades da implementação ROS 1. A partir dessas relações, foram caracterizadas três transformações identificadas na implementação ROS 2: a adoção de mecanismos de gerenciamento do ciclo de vida dos nós, a ampliação das responsabilidades de monitoramento e a reorganização das operações de processamento, comunicação e avaliação de risco entre componentes especializados. Essa análise permitiu examinar como responsabilidades associadas aos comportamentos selecionados foram mantidas, adaptadas ou redistribuídas na plataforma de destino. Complementarmente, a comparação dos cenários BDD, de seus oráculos e das evidências de execução permitiu identificar correspondências e discrepâncias comportamentais, delimitando o alcance das verificações realizadas. Os resultados demonstram a possibilidade de integrar essas evidências na análise retrospectiva da migração, sem estabelecer que a abordagem tenha orientado as decisões originais ou que sua eficácia tenha sido demonstrada em outros sistemas.




